% resume.tex
%
% (c) 2002 Matthew Boedicker <mboedick@mboedick.org> (original author) http://mboedick.org
% (c) 2003-2007 David J. Grant <davidgrant-at-gmail.com> http://www.davidgrant.ca
%
% This work is licensed under the Creative Commons Attribution-ShareAlike 3.0 Unported License. To view a copy of this license, visit http://creativecommons.org/licenses/by-sa/3.0/ or send a letter to Creative Commons, 171 Second Street, Suite 300, San Francisco, California, 94105, USA.

\documentclass[letterpaper,11pt]{article}

\newlength{\outerbordwidth}
\pagestyle{empty}
\raggedbottom
\raggedright
\usepackage[svgnames]{xcolor}
\usepackage{framed}
\usepackage{tocloft}
\usepackage{amssymb}
\usepackage{etoolbox}
\usepackage{hyperref, url}
\robustify\cftdotfill

%-----------------------------------------------------------
%Edit these values as you see fit
\setlength{\outerbordwidth}{3pt}  % Width of border outside of title bars
\definecolor{shadecolor}{gray}{0.75}  % Outer background color of title bars (0 = black, 1 = white)
\definecolor{shadecolorB}{gray}{0.93}  % Inner background color of title bars

%-----------------------------------------------------------
%Margin setup
\setlength{\evensidemargin}{-0.25in}
\setlength{\headheight}{-0.25in}
\setlength{\headsep}{0in}
\setlength{\oddsidemargin}{-0.25in}
\setlength{\paperheight}{11in}
\setlength{\paperwidth}{8.5in}
\setlength{\tabcolsep}{0in}
\setlength{\textheight}{9.75in}
\setlength{\textwidth}{7in}
\setlength{\topmargin}{-0.3in}
\setlength{\topskip}{0in}
\setlength{\voffset}{0.1in}

%-----------------------------------------------------------
%Custom commands
\newcommand{\resitem}[1]{\item #1 \vspace{-2pt}}
\newcommand{\resheading}[1]{\vspace{8pt}
  \parbox{\textwidth}{\setlength{\FrameSep}{\outerbordwidth}
    \begin{shaded}

\setlength{\fboxsep}{0pt}\framebox[\textwidth][l]{\setlength{\fboxsep}{4pt}\fcolorbox{shadecolorB}{shadecolorB}{\textbf{\sffamily{\mbox{~}\makebox[6.762in][l]{\large #1} \vphantom{p\^{E}}}}}}
    \end{shaded}
  }\vspace{-5pt}
}

\newcommand{\ressubheading}[4]{
\begin{tabular*}{6.5in}{l@{\cftdotfill{\cftsecdotsep}\extracolsep{\fill}}r}
		\textbf{#1} & #2 \\
		\textit{#3} & \textit{#4} \\
\end{tabular*}\vspace{-6pt}}

\begin{document}

\begin{tabular*}{7in}{l@{\extracolsep{\fill}}r}
\textbf{\Large Abhijit Chowdhary} & (240)-715-8308\\
Department of Mathematics & achowdh2@ncsu.edu \\
North Carolina State University & \url{https://abhijit-c.github.io/}\\
Raleigh, NC & 
\end{tabular*}
\\

\vspace{0.1in}

\resheading{Education}
\begin{itemize}
\item
	\ressubheading{North Carolina State University}
  {Raleigh, NC}
  {Applied Math PhD (GPA: 4.0)}
  {Aug. 2020 - Present}
  \begin{itemize}
    \resitem{
      (Aug. 2021 - Present) --- RA for
      \href{https://aalexan3.math.ncsu.edu/}{Alen Alexanderian} 
      via NSF-DMS-2111044
    }
    \resitem{
      (Aug. 2021) --- Passed Qualification Exams (Matrix Theory, PDE, Analysis)
    }
    \resitem{(Aug. 2020 - May 2021) --- Recipient of the Siewert Fellowship}
    \resitem{
      Coursework: Analysis I/II, PDE I/II, Matrix Theory I/II, Modeling
      I/II, Probability I, Uncertainty Quantification, Computational Inverse
      Problems, Advanced Functional Analysis, Numerical Nonlinear PDEs,
      Data-Driven Modeling and Analysis
    }
  \end{itemize}

\item
	\ressubheading{New York University}{New York, NY}{B.A. Joint Mathematics and
    Computer Science, Classics Minor}{Sep.2016 - May.2020}
	\begin{itemize} 
    \resitem{(In Major GPA: 3.657)}
		\resitem{Relevant undergraduate courses: 
      Algorithms, 
      Chaos and Dynamical Systems,
      Computer System Organization, 
      Operating Systems,
      Honors Algebra I/II, 
      Honors Analysis I/II, 
      Honors Linear Algebra, 
      Honors Probability Theory, 
      Numerical Computation, 
      Topology.}
    \resitem{Relevant graduate courses: 
      Algebra,
      Basic Probability,
      Convex and Nonsmooth Optimization,
      Finite Element Method,
      Fundamental Algorithms, 
      Geometric Modeling, 
      High Performance Computing,
      Methods of Applied Math, 
      Numerical Methods I/II, 
      Partial Differential Equations.}
	\end{itemize}

\item
	\ressubheading{University of Maryland}{College Park, MD}{Visiting
    Student (GPA: 3.925)}{Summer 2017,2018}
	\begin{itemize}
		\resitem{Coursework: Complex analysis, Number theory, Partial
        Differential Equations, Introduction to Artificial Intelligence.}
	\end{itemize}

\end{itemize}

\resheading{Research Interests}
High performance computing and parallel scientific computing;
Numerical solution of ODEs and PDEs;
PDE-constrained optimization;
Inverse problems;
Uncertainty quantification;
Randomized numerical methods;
Machine learning for scientific applications

\resheading{Research Activity}
\begin{itemize}
  \item 
  \ressubheading{Robust Optimal Experimental Design for Non-Linear Bayesian Inference}{}{(Argonne National Laboratory) Research with Ahmed Attia}{May 2023 - Present}
  \begin{itemize}
    \resitem{
	During my Givens internship at Argonne National Lab during the summer of 2023, I began joint work with Dr. Ahmed Attia on developing scalable algorithms for solving optimal experimental design problems for non-linear bayesian inversion.
    }
  \resitem{
	This includes work and continued development on the open source package \href{https://gitlab.com/ahmedattia/pyoed/}{PyOED}. This package is to be presented at a mini-symposia in UQ24.
  }
  \end{itemize}
  \item 
  \ressubheading{HDSA enabling the Kullback Leibler Divergence as an UQ
  objective}{}{(NCSU) Research with Alen Alexanderian}{Aug. 2021 - Present}
  \begin{itemize}
    \resitem{
      As a main objective as a research assistant for Dr. Alen Alexanderian,
      I worked on implementing theory from hyper-differential sensitvity
      analysis to enable taking efficient and exact derivaitves via
      adjoint-based gradient computation of the KL-Divergence in order to
      undertand its sensitivity in the context of uncertainty quantification.
    }
    \resitem{
      In the process, using FENiCS (Python) I implemented a Bayesian linear
      inverse problem for a 3D linear elasticity model, stemming from the
      real-world problem of earthquake modeling, that served as
      a proof-of-concept of the theory. This model is planned to be reused in
      further work in Bayesian inverse problems as an interesting test problem.
    }
    \resitem{
      Paper with the resulting work is currently in writing.
    }
  \end{itemize}

\item
	\ressubheading{Math REU: Imperfect Periodic Patterns}{Athens, OH}{Ohio
  University}{June 2019 - August 2019}
	\begin{itemize}
    \resitem{I joined a research team under professor Qiliang Wu, and another
      undergraduate Mason Haberle from Berkeley in researching the field of
      pattern formation.}
    \resitem{Our team specifically set out to prove nonlinear stability of the
      2D Swift-Hohenberg equation at the zigzag boundary, and as of now we've
      completed the proof and the paper is in the draft stages.}
    \resitem{The challenge here was mostly on how to adapt known techniques
      first to the Swift-Hohenberg equation, and second to higher dimensions.
      This was mostly a conceptual difficulty in the functional analysis
      framework surrounding the current research which we had to resolve.}
    \resitem{\href{https://arxiv.org/abs/2012.04154}{The report can be found here.}}
	\end{itemize}
\end{itemize}

\resheading{Presentations}
\begin{itemize}
  \item \ressubheading{Robust Optimal Experimental Design for Non-Linear Robust
  OED}{Argonne National Lab}{SASSy Seminar}{July 28th 2023}
  \item \ressubheading{Sensitivity Analysis of the Info. Gain in
  Inf. Dim. Bayesian Inference}{Amsterdam}{SIAM CSE23}{Mar. 2nd 2023}
  \item \ressubheading{Eigenvalue Sensitivities for the Info. Gain in
  Linear Bayesian Inference}{Raleigh, NC}{Applied Math Graduate Student
  Seminar}{Sep. 26th 2022}
  \item \ressubheading{Inf. Dim. Bayesian Inference for Fault Slip
  from Surface Measurements}{Raleigh, NC}{Applied Math Graduate Student
  Seminar}{Apr. 25th 2022}
\end{itemize}

\resheading{Projects}
\begin{itemize}

\item
  \ressubheading{Benchmarking Reservoir Computing for Lorenz-96}{}{(NCSU)
  Numerical Nonlinear PDEs Course}{Nov. 2022 - Dec. 2022}
  \begin{itemize}
    \resitem{
      For a class final project, I investigate the performance of Reservoir
      computing, a simplified method of learning derived from traditional
      recurrent neural network (RNN) ideas, on the Lorenz-96 toy model.
    }
    \resitem{
      As the Lorenz-96 model is a model with comparable dynamics to more complex
      athmospheric models, the goal of this project was to demonstrate time
      series forcasting prowess of of the method, reminiscint of traditional
      RNNs while also demonstrating how much quicker the model was to train.
    }
    \resitem{\href{https://abhijit-c.github.io/files/ddm-final.pdf}{Link
    to the report.}}
  \end{itemize}

\item
  \ressubheading{Well-Balanced Stochastic Galerkin for random PDEs}{}{(NCSU)
  Numerical Nonlinear PDEs Course}{Apr. 2022 - May 2022}
  \begin{itemize}
    \resitem{
      For a class final project, I along with Andrew Shedlock (NCSU) worked on
      an implementation of a well-balanced stochastic Galerkin method for PDEs
      with random forcing.
    }
    \resitem{
        Verified the results in the original paper, and did smaller studies on
        computational efficiency and noted places in the algorithm where one
        could take advantage of symmetry to improve runtime.
    }
    \resitem{
      Implemented first in MATLAB, and later rewritten in Julia to avoid
      unecessary allocations and take advantage of its pre-compilation step,
      achieving a 15x speedup.
    }
    \resitem{\href{https://abhijit-c.github.io/files/WBSG-Report.pdf}{Link
    to the report.}}
  \end{itemize}

\item
  \ressubheading{HOGWILD!}{}{(NYU) Convex and Nonsmooth Optimization}{Apr. 2020 - May 2020}
	\begin{itemize}
    \resitem{
      For a class final project, I looked into a lock-free parallel stochastic
      gradient descent implementation, {\rm HOGWILD!}.
    }
    \resitem{
      Implemented the algorithm, and did a case study on the convergence and
      efficiency analysis. The goal was the to try and study the induced
      asynchronous noise as a function of bandedness in order to apply them to
      banded matrices arising from a discretization of a Poisson equation, but
      COVID-19 cut the project short.
    }
    \resitem{Implementation with Eigen and OpenMP in C++.}
    \resitem{\href{https://abhijit-c.github.io/files/Hogwild.pdf}{Link
    to the report.}}
	\end{itemize}



\item
  \ressubheading{Parareal}{}{(NYU) High Performance Computing and Numerical Methods
  II}{Apr. 2019 - May 2019}
	\begin{itemize}
    \resitem{For a class final project, I decided to look into parallel
    techniques for solving ordinary differential equations, in particular the
    parallel-in-time algorithm, \textit{Parareal}.}
    \resitem{For this project, I implemented and analyzed this algorithm, and
    further tested it's scaling properties on the HPC cluster Prince here at
    NYU.}
    \resitem{Implemented with Eigen and OpenMP in C++}
    \resitem{\href{https://abhijit-c.github.io/files/Parareal/Parareal.pdf}{Link
    to the report.}}
	\end{itemize}

\item
	\ressubheading{Algebraic Point Set Surfaces Implementation}{}{(NYU) Geometric
    Modeling}{Apr. 2018 - May. 2018}
	\begin{itemize}
		\resitem{For a class final project, I implemented the theory in the
        paper \textit{Algebraic Point Set Surfaces} by Ga\={e}l Gunnebaud and
        Markus Gross from ETH Zurich.}
    \resitem{The Paper presented an alternative method to take a point cloud
      to a triangularized mesh, and another method to estimate normals from a
      point cloud using algebraic fitting of a sphere.}
		\resitem{This was mostly a challenge in comprehension of the paper and
      implementation, notably fighting with Eigen to try and constuct and
      solve the sytems in an efficient manner.}
    \resitem{Implemented with Eigen, libigl, and nanoflann in C++}
	\end{itemize}

\end{itemize}

\resheading{Teaching and Service}

\begin{itemize}

\item 
	\ressubheading{Organizing the Applied Math Student Seminar (AMGSS)}{Raleigh, NC}{NCSU}{Fall 2022 - Present}

\item
	\ressubheading{Recitation Leader and Grader for Calculus II}{Raleigh, NC}{NCSU MA 141}{Spring and Summer 2021}

\item
	\ressubheading{Recitation Leader and Grader for Calculus}{Raleigh, NC}{NCSU MA 121}{Fall 2020}

\item
	\ressubheading{Tutor for Numerical Computing}{New York, NY}{NYU CSCI-UA 421}{Spring 2020}

\item
	\ressubheading{Tutor and TA for Algorithms}{New York, NY}{NYU, CSCI-UA 310,
	CSCI-GA 1170 }{Sep. 2017 - May 2019}
	\begin{itemize}
		\resitem{
      Worked as a Tutor and TA to Professor Siegel at NYU for his undergraduate
      basic algorithms and graduate fundamental algorithms course.
    }
	\end{itemize}

\item 
	\ressubheading{EBoard of ACM@NYU}{New York, NY}{NYU}{Fall 2017 - Spring 2020}

\item
	\ressubheading{First Robotics Team Member, Team 2849: Ursa Major}{Columbia, MD}{Hammond High
    School}{Sep. 2012 - 2019}
	\begin{itemize}
		\resitem{A robotics team; every new year they gather for a challenge
        created by FIRST Robotics to build a robot in six weeks.}
        \resitem{I worked as a build-team / programming-team flex member and
        team captain during my student years, and now I help as a programming
        and design mentor during their season.}
        \resitem{Has managed to consistantly reach eliminantion and championship
        rounds at the regional level.}
        \resitem{\href{https://github.com/teamursamajor}{Link to their Github}}
	\end{itemize}
\end{itemize}

\resheading{Skills}

\begin{description}
\item[Programming Languages:]
C, C++, Julia, Matlab, Python and a little Mathematica.
\item[Libraries and Technologies:]
  Num/Scipy, FEniCS, OpenMP, MPI, HPC Tooling (slurm, etc.), CUDA, Linux (have
  been using for over 8 years, currently on Fedora).
\item[Languages:]
English, Latin, Broken spoken Hindi
\item[Minor Mechanical Fabrication Skills]
\end{description}

\end{document}
